\subsection{Shift magnitude}
\label{sec:magnitude}

Figure~\ref{fig:magnitude} varies the strength of each shift. Thresholds are matched for the comparators and at the theoretical value for TSSD; each curve is the fraction of 200 streams detected (for TSSD and the ensemble, detected \emph{and} correctly typed).

\emph{Covariate.} TSSD is the most sensitive covariate detector in the study. At one tenth of the base translation it detects and correctly types 114/200 streams, against 44 for the best specialist (kernel two-sample betting) and 3 for the ensemble. At 0.2 it types 200/200. The ensemble's typed accuracy stays poor at every magnitude (79/200 at half strength), because prevalence-driven CUSUM alarms label most covariate shifts as label shifts.

\emph{Label.} Mild prevalence changes are TSSD's weakest regime. At $\gamma=0.5$ the classifier-CUSUM detects 188/200 and TSSD 59/200, and only at $\gamma=0.75$ does TSSD reach 182/200. The label witness has to work through the gate and the orthogonalised covariate witness, while the CUSUM is optimised for exactly this alternative.

\emph{Partial concept.} When only a fraction of labels is cycled, the certified concept claim needs much more evidence than an uncorrected loss monitor. With half of the labels cycled, TSSD certifies 76/200 streams while Page--Hinkley detects 194/200; at 10\% cycling TSSD certifies essentially none. This is the clearest instance of the price of certification. A partial label permutation lies close to the prior-correctable family, and with about ten post-change labels the null fitted with hindsight absorbs most of the evidence.

\begin{figure}[t]
\centering
\begin{tikzpicture}
\begin{groupplot}[group style={group size=3 by 1, horizontal sep=0.9cm},
  width=0.36\textwidth, height=0.3\textwidth, ymin=-0.02, ymax=1.02,
  ylabel style={font=\small}, xlabel style={font=\small}, tick label style={font=\scriptsize},
  title style={font=\small}, legend style={font=\scriptsize, at={(0.5,-0.42)}, anchor=north, legend columns=5, draw=none},
  every axis plot/.append style={thick, mark size=1.6pt}]
\nextgroupplot[title={Covariate}, xlabel={translation scale $s$}, ylabel={fraction of streams}, xmode=log, xtick={0.1,0.2,0.5,1}, xticklabels={0.1,0.2,0.5,1}]
\addplot[black, mark=*] table[x=level,y=TSSDSR]{figures/mag_covariate.dat};
\addplot[black, dashed, mark=o] table[x=level,y=Ensemble]{figures/mag_covariate.dat};
\addplot[blue, mark=square*] table[x=level,y=Kernel2ST]{figures/mag_covariate.dat};
\addplot[orange, mark=triangle*] table[x=level,y=BBSD]{figures/mag_covariate.dat};
\addplot[red!70!black, mark=diamond*] table[x=level,y=CUSUM]{figures/mag_covariate.dat};
\nextgroupplot[title={Label}, xlabel={prior mixing $\gamma$}]
\addplot[black, mark=*] table[x=level,y=TSSDSR]{figures/mag_label.dat};
\addplot[black, dashed, mark=o] table[x=level,y=Ensemble]{figures/mag_label.dat};
\addplot[blue, mark=square*] table[x=level,y=Kernel2ST]{figures/mag_label.dat};
\addplot[orange, mark=triangle*] table[x=level,y=BBSD]{figures/mag_label.dat};
\addplot[red!70!black, mark=diamond*] table[x=level,y=CUSUM]{figures/mag_label.dat};
\nextgroupplot[title={Concept}, xlabel={fraction of labels cycled}]
\addplot[black, mark=*] table[x=level,y=TSSDSR]{figures/mag_concept.dat};
\addplot[black, dashed, mark=o] table[x=level,y=Ensemble]{figures/mag_concept.dat};
\addplot[green!50!black, mark=pentagon*] table[x=level,y=PH]{figures/mag_concept.dat};
\addplot[violet, mark=x] table[x=level,y=CTM]{figures/mag_concept.dat};
\addplot[cyan!70!black, mark=star] table[x=level,y=WATCH]{figures/mag_concept.dat};
\end{groupplot}
\node[font=\scriptsize, anchor=north, align=center] at ($(group c2r1.south)+(0,-1.05cm)$) {%
\tikz{\draw[black,thick](0,0)--(0.4,0);\fill (0.2,0) circle (1.6pt);} TSSD (SR, typed)\quad
\tikz{\draw[black,thick,dashed](0,0)--(0.4,0);\draw (0.2,0) circle (1.6pt);} specialist ensemble (typed)\quad
\tikz{\draw[blue,thick](0,0)--(0.4,0);} kernel 2ST\quad
\tikz{\draw[orange,thick](0,0)--(0.4,0);} BBSD\\
\tikz{\draw[red!70!black,thick](0,0)--(0.4,0);} CUSUM\quad
\tikz{\draw[green!50!black,thick](0,0)--(0.4,0);} Page--Hinkley\quad
\tikz{\draw[violet,thick](0,0)--(0.4,0);} conformal TM\quad
\tikz{\draw[cyan!70!black,thick](0,0)--(0.4,0);} WATCH-style};
\end{tikzpicture}

\caption{Detection versus shift magnitude (Gaussian, 200 streams per point; level 1 is the main study). TSSD and the ensemble count typed-correct detections; the other detectors count detections. Comparators use matched 5\% thresholds; TSSD uses its theoretical SR threshold.}
\label{fig:magnitude}
\end{figure}

\subsection{Label budget}
\label{sec:reveal}

Figure~\ref{fig:reveal}(a) and Table~\ref{tab:reveal} vary the reveal rate; Figure~\ref{fig:reveal}(b) shows the three test-mode evidence processes on one label+concept stream. The covariate and label witnesses hardly depend on it. The certified concept claim depends on it strongly: it types 70/200 concept streams correctly at 2\%, 180/200 at 5\%, 198/200 at 10\% and 199/200 at 20\%, where the median delay drops to 22 steps. From 10\% upwards it matches the conformal test martingale's detection rate, while also certifying the type. At 2\% the clipped paired witness of the ablation detects only 5/200. Across all four rates, no concept claim was made on any of the 4400 streams whose posterior is prior-correctable.

\begin{figure}[t]
\centering
\begin{tikzpicture}
\begin{groupplot}[group style={group size=2 by 1, horizontal sep=1.6cm},
  width=0.46\textwidth, height=0.32\textwidth, tick label style={font=\scriptsize},
  xlabel style={font=\small}, ylabel style={font=\small}, title style={font=\small},
  every axis plot/.append style={thick, mark size=1.6pt}]
\nextgroupplot[title={(a) Concept arms vs.\ label rate}, xlabel={label reveal rate $\rho$}, ylabel={fraction of 200 streams},
  xmode=log, xtick={0.02,0.05,0.1,0.2}, xticklabels={2\%,5\%,10\%,20\%}, ymin=-0.02, ymax=1.02,
  legend style={font=\scriptsize, at={(0.98,0.03)}, anchor=south east, draw=none, fill=white, fill opacity=0.8, text opacity=1}]
\addplot[black, mark=*] table[x=rate,y=con_s]{figures/reveal.dat}; \addlegendentry{TSSD SR (typed)}
\addplot[black, dashed, mark=o] table[x=rate,y=con_t]{figures/reveal.dat}; \addlegendentry{TSSD test (typed)}
\addplot[gray, mark=square] table[x=rate,y=con_p]{figures/reveal.dat}; \addlegendentry{paired $C$ ablation}
\addplot[violet, mark=x] table[x=rate,y=con_ctm]{figures/reveal.dat}; \addlegendentry{conformal TM (untyped)}
\nextgroupplot[title={(b) Example label+concept stream}, xlabel={step $t$}, ylabel={log evidence (test mode)},
  ymin=-2, ymax=30, xmin=0, xmax=400,
  legend style={font=\scriptsize, at={(0.02,0.98)}, anchor=north west, draw=none, fill=white, fill opacity=0.8, text opacity=1}]
\addplot[blue, no marks] table[x=t,y=X]{figures/trace.dat}; \addlegendentry{$X$}
\addplot[orange, no marks] table[x=t,y=L]{figures/trace.dat}; \addlegendentry{$L$}
\addplot[red!70!black, no marks, very thick] table[x=t,y=C]{figures/trace.dat}; \addlegendentry{$C$ (certified)}
\addplot[black, dotted, no marks, domain=0:400] {ln(60)}; \addlegendentry{$\log(3/\alpha)$}
\draw[gray, dashed] (axis cs:200,-2) -- (axis cs:200,30);
\end{groupplot}
\end{tikzpicture}

\caption{(a) Typed-correct concept detections versus label budget (Gaussian). (b) Test-mode log evidence on one label+concept stream; the change is at step 201 (dashed). The certified concept witness crosses about 70 steps after the change and fixes the reported action type as concept; the label witness crosses only near the horizon, and the covariate witness stays flat.}
\label{fig:reveal}
\end{figure}

\begin{table}[t]
\centering
\caption{Label-budget sensitivity (Gaussian, 200 streams per cell; theory thresholds). Entries are typed-correct detections; ``null'' is the number of null streams with any alarm. CTM = conformal test martingale (native threshold), counted as detections.}
\label{tab:reveal}
\small
\begin{tabular}{@{}llrrrr@{}}
\toprule
Rate & Condition & TSSD SR & TSSD test & Paired $C$ (ablation) & CTM \\
\midrule
\tabinput{results/tables/reveal_rows.tex}
\bottomrule
\end{tabular}
\end{table}

\subsection{Change time: test mode versus SR mode}
\label{sec:changetime}

Table~\ref{tab:changetime} moves the change point from step 100 to step 2000 (100 streams per cell, horizon $t_0+300$). The test-mode delay grows with the length of the null period, from 35 to 147 steps for covariate shift and from 108 to 186 for label shift. The test martingale's wealth decays under the null, so a late change finds it far below its starting point. The SR delay is flat: about 23 steps for covariate shift and about 75 for label shift, whatever the change time.

The concept witness is almost unaffected in either mode. Before the change its alternative tracks the null closely, so its evidence barely decays.

The price of restarting is a higher alarm rate over long null stretches. On 200 null streams of 2300 steps, SR mode raises 23 alarms and test mode 8. Both are consistent with their guarantees: an ARL of at least $m/\alpha=8000$ steps for SR, and at most 5\% ever for test mode. Within the 400-step design window each mode raises 3. These pre-change alarms are the main reason SR detection falls from 99--100 to 90--94 streams at $t_0=2000$. For long deployments, $A$ can be set from the desired ARL rather than from a fixed window. Across all 1000 change-time streams with a prior-correctable posterior, no concept claim was made.

\begin{table}[t]
\centering
\caption{Change time (Gaussian, 100 streams per cell). Typed-correct detections with median delay in parentheses, test mode versus SR mode.}
\label{tab:changetime}
\small
\begin{tabular}{@{}lrrrrrr@{}}
\toprule
 & \multicolumn{2}{c}{Covariate} & \multicolumn{2}{c}{Label} & \multicolumn{2}{c}{Concept} \\
\cmidrule(lr){2-3}\cmidrule(lr){4-5}\cmidrule(lr){6-7}
$t_0$ & test & SR & test & SR & test & SR \\
\midrule
\tabinput{results/tables/changetime_rows.tex}
\bottomrule
\end{tabular}
\end{table}



\subsection{Ablations}
\label{sec:ablation}

Table~\ref{tab:ablation} perturbs each design constant on fresh streams: 300 nulls and 100 streams per shift. The default configuration corresponds to the main study (Table~\ref{tab:typing}).

No perturbation produced a false concept claim. Null alarm rates stayed at or below $4/300$. The covariate, label and covariate+label results move by at most a few streams in SR mode, and the test-mode label witness is sensitive only to a tighter gate (65 versus 86 typed-correct).

The concept results reveal a robustness--power trade-off in the null family. Removing tempering ($B=\{1\}$) or narrowing the prior box ($|u_k|\le4$) raises the typed concept rate from 87 to 97--99. A wider tempering grid lowers it to 76. On the well-specified Gaussian model the smaller family is safe. We nevertheless keep the tempered family as the default, because it is what protects the concept claim against moderate miscalibration on real data, where the source model is not well specified. The two-component alternative without the refit expert performs similarly on Gaussian data (85 versus 87); the refit expert matters most when the posterior change involves directions that the frozen logits do not represent.

\begin{table}[t]
\centering
\caption{Ablations (Gaussian; 300 nulls, 100 streams per shift). Each cell gives SR / test-mode values: null streams with any alarm, or typed-correct detections. The last column counts concept claims (either mode) on the 700 streams whose posterior is prior-correctable.}
\label{tab:ablation}
\scriptsize
\setlength{\tabcolsep}{3.5pt}
\begin{tabular}{@{}lrrrrrrr@{}}
\toprule
Perturbation & Null & Cov & Lab & Con & C+L & L+C & False $C$ \\
\midrule
\tabinput{results/tables/ablation_rows.tex}
\bottomrule
\end{tabular}
\end{table}


\subsection{Real data}
\label{sec:real}

\paragraph{Digits.} Table~\ref{tab:digits} reports the exact-validity Digits study. The null alarm rates are 1/300 (SR) and 3/300 (test) at 5\% labels, and 2/100 and 3/100 at 20\%. The acquisition corruption is detected and typed as covariate in all streams. Prevalence change is TSSD's weak spot here: SR mode types 53/100 label streams correctly against 99/100 detections by classifier-CUSUM. The test mode is weaker still, because the covariate gate often closes before the label evidence accumulates. With ten classes and about seven post-change labels at 5\%, the certified concept claim fires on only 12--31\% of the concept alternatives. At 20\% labels it fires on 100/100 in each of the three concept arms, and it makes no false claim on any null, covariate, label or covariate+label stream at either rate. Unlike on the Gaussian data, the specialist ensemble types Digits well at 5\% labels: CUSUM is silent under the acquisition corruption, so the ensemble's characteristic failure does not arise. Its concept reports remain uncertified, however\footnote{The ensemble's thresholds are matched at 5\% labels. At 20\% its null alarm rate rises to 32/100, so its 20\% entries are not shown.}.

\begin{table}[t]
\centering
\caption{Digits (exact-validity protocol, 100 streams per shift). Entries are typed-correct detections for TSSD and the ensemble and detections for the other methods. Comparators use matched thresholds except CTM and WATCH-style, which use their native anytime-valid thresholds.}
\label{tab:digits}
\small
\setlength{\tabcolsep}{4pt}
\begin{tabular}{@{}lrrrrrr@{}}
\toprule
 & Cov & Lab & Con & C+L & C+C & L+C \\
\midrule
\multicolumn{7}{@{}l}{\emph{5\% labels}}\\
\tabinput{results/tables/digits_rows05.tex}
\midrule
\multicolumn{7}{@{}l}{\emph{20\% labels (theory/native thresholds)}}\\
\tabinput{results/tables/digits_rows20.tex}
\bottomrule
\end{tabular}
\end{table}

\paragraph{Covertype: natural geographic shift.} On exact same-domain nulls (held-out Rawah data), TSSD alarms on 3/300 (SR) and 2/300 (test) streams. The ensemble, calibrated on disjoint nulls of the same kind, alarms on 27/300. All three natural target domains are detected by TSSD in 96--100\% of streams (median SR delay 31--57 steps). They are reported as covariate or label shifts, and the certified concept claim is never made at 5\% labels and is made in 3--4 of 30 streams at 20\%.

The offline oracle (Table~\ref{tab:oracle}) explains this. In every domain the refitted model beats the best prior correction by 0.19--0.37 nats per case, so the oracle type is concept. But a prior correction already removes 43\%, 75\% and 98\% of the frozen model's excess loss in areas 1, 2 and 3. Certifying the remaining gap of a few tenths of a nat per label, against a null fitted with hindsight over seven classes, needs more labels than the 20--80 available. The ensemble reports concept on most area-2 and area-3 streams. This agrees with the oracle, but the ensemble's concept claims are uncertified: the same rule produced 149 false concept claims on the Gaussian streams. We read Covertype as an honest limitation. TSSD's ``covariate'' state means \emph{inputs moved and no posterior change is certified}, not \emph{the model is safe}. When harm matters more than type, it should be combined with a risk monitor (Section~\ref{sec:discuss}).

\begin{table}[t]
\centering
\caption{Offline full-label characterisation of the natural shifts (not available to any monitor). LL: log loss on held-out target data (nats per case); prior: best prior and tempering correction in the family \eqref{eq:family}; refit: logistic model refitted on target data; AUC: input domain classifier.}
\label{tab:oracle}
\small
\begin{tabular}{@{}lrrrrrrl@{}}
\toprule
Domain / segment & acc.\ frozen & LL frozen & LL prior & LL refit & gap & AUC & oracle \\
\midrule
\tabinput{results/tables/oracle_rows.tex}
\bottomrule
\end{tabular}
\end{table}

\paragraph{Electricity: chronological drift.} Table~\ref{tab:elec} monitors 18 consecutive segments of the Elec2 market data in time order.

On exchangeable source-period nulls, TSSD alarms on 2/300 (SR) and 4/300 (test) streams. All comparators were calibrated to 5\% on these nulls.

The offline oracle splits the segments into two regimes. In the first five (records 8000--17999) the input distribution has moved (AUC 0.55--0.79) and a prior correction removes most of the excess loss (gap $\le0.053$ nats). In the remaining 13 (records 18000 onwards) the frozen posterior fails badly (frozen log loss 5.7--15.3) and only a refit repairs it (gap 2.2--8.8 nats).

TSSD follows this structure exactly. In the first regime it raises an alarm in every run and reports a covariate shift, never a concept shift. In the second regime it makes a \emph{certified} concept claim in all 260 runs (median SR delay 24 records). The loss-based detectors agree on the second regime: Page--Hinkley fires in all of its runs, and the conformal martingale in 15--20 of 20 per segment. The WATCH-style weighted martingale fires in only 4--9 of 20 runs on four of these segments. Its covariate weighting absorbs part of the change when the inputs move strongly, the behaviour it is designed to have for benign covariate shift. Among the input detectors, BBSD fires in both regimes. TSSD is the only method whose output distinguishes them with a guarantee.

The early segments use labels that are serially dependent and only roughly exchangeable, so the absence of concept claims there is descriptive rather than a test of validity.

\begin{table}[t]
\centering
\caption{Electricity: chronological segments of 1000 records (20 runs each, 5\% labels). The oracle type and gap come from Table~\ref{tab:oracle}'s procedure. TSSD columns give alarm counts and the SR action-type counts concept/label/covariate. Comparator columns give alarm counts at matched thresholds (PH = Page--Hinkley; CTM = conformal TM; WATCH = WATCH-style WCTM).}
\label{tab:elec}
\scriptsize
\setlength{\tabcolsep}{3.5pt}
\begin{tabular}{@{}lrrrrcrrrr@{}}
\toprule
Segment start & Oracle & Gap & SR alarm & Test alarm & C / L / X & PH & CTM & WATCH & BBSD \\
\midrule
\tabinput{results/tables/elec_rows.tex}
\bottomrule
\end{tabular}
\end{table}


